
\documentclass[../../../physics-standard_questions_all.tex]{subfiles}


\begin{document}

図のように，真空中において，同形の4枚の金属極板X，Y，Z，Wがあり，
XとYが間隔$2d$で，ZとWが間隔$d$で平行に向き合わせてある．
それらは，電池B（起電力$V_0$），スイッチS$_1$，S$_2$，
およびアースと接続されて回路を形成している．
極板の面積を$S$，真空の誘電率を$\varepsilon_0$とし，端の効果は無視する．
はじめ，S$_1$，S$_2$は共に開いており，いずれの極板も帯電していない．

\begin{enumerate}

    \item S$_1$を閉じて安定するまで待った．
    このとき，X下面の電荷$x_1$，Z下面の電荷$z_1$を求めよ．
    また，XY間の電場$E_1$（X→Yの向き正）を求めよ．
    
    \item S$_1$を開き，S$_2$を閉じて安定するまで待った．
    このとき，X下面の電荷$x_2$，Z下面の電荷$z_2$を求めよ．
    また，XY間の電場$E_2$（X→Yの向き正），Xの電位$\phi_2$を求めよ．

    \item S$_2$を開き，S$_1$を閉じて安定するまで待った．
    このとき，X下面の電荷$x_3$，Z下面の電荷$z_3$を求めよ．
    
    \item S$_1$を閉じたまま，S$_2$も閉じて安定するまで待った．
    このとき，X下面の電荷$x_4$，Z下面の電荷$z_4$を求めよ．

    \item S$_2$を閉じたまま，S$_1$を開き，ゆっくりとZを移動させ，
    ZとWの間隔が$2d$になるようにした．
    このとき，X下面の電荷$x_5$，Z下面の電荷$z_5$を求めよ．

\end{enumerate}

\vspace{0.2\baselineskip}
{\centering
    \includegraphics{\subfix{fig/fig_em_cb_06/fig_em_cb_06_00.pdf}}
\par}
\vspace{0.2\baselineskip}

\end{document}



